\documentclass{beamer}

\usepackage[russian]{babel}
\usepackage[utf8]{inputenc}
\usepackage{cmap}
\usepackage{graphicx}
\usepackage{xspace}
\usepackage{psfrag}
\usepackage{hyperref}
\usepackage[normalem]{ulem}

\beamertemplatenavigationsymbolsempty
\setbeamertemplate{footline}[frame number]
\usecolortheme{seahorse}

%%Исправляем греческие буквы в формулах, они должны быть прямыми, а не курсивными
\DeclareSymbolFont{greek}{U}{eur}{m}{n}                                                        
\SetSymbolFont{greek}{bold}{U}{eur}{b}{n}                                                      
\DeclareSymbolFontAlphabet{\gr}{greek}                                                         
\DeclareMathSymbol{\alpha}{\mathord}{greek}{"0B}                                               
\DeclareMathSymbol{\beta}{\mathord}{greek}{"0C}                                                
\DeclareMathSymbol{\gamma}{\mathord}{greek}{"0D}                                               
\DeclareMathSymbol{\delta}{\mathord}{greek}{"0E}                                               
\DeclareMathSymbol{\epsilon}{\mathord}{greek}{"0F}                                             
\DeclareMathSymbol{\zeta}{\mathord}{greek}{"10}                                                
\DeclareMathSymbol{\eta}{\mathord}{greek}{"11}                                                 
\DeclareMathSymbol{\theta}{\mathord}{greek}{"12}                                               
\DeclareMathSymbol{\iota}{\mathord}{greek}{"13}                                                
\DeclareMathSymbol{\kappa}{\mathord}{greek}{"14}                                               
\DeclareMathSymbol{\lambda}{\mathord}{greek}{"15}                                              
\DeclareMathSymbol{\mu}{\mathord}{greek}{"16}                                                  
\DeclareMathSymbol{\nu}{\mathord}{greek}{"17}                                                  
\DeclareMathSymbol{\xi}{\mathord}{greek}{"18}                                                  
\DeclareMathSymbol{\pi}{\mathord}{greek}{"19}                                                  
\DeclareMathSymbol{\rho}{\mathord}{greek}{"1A}                                                 
\DeclareMathSymbol{\sigma}{\mathord}{greek}{"1B}                                               
\DeclareMathSymbol{\tau}{\mathord}{greek}{"1C}                                                 
\DeclareMathSymbol{\upsilon}{\mathord}{greek}{"1D}                                             
\DeclareMathSymbol{\phi}{\mathord}{greek}{"1E}                                                 
\DeclareMathSymbol{\chi}{\mathord}{greek}{"1F}                                                 
\DeclareMathSymbol{\psi}{\mathord}{greek}{"20}                                                 
\DeclareMathSymbol{\omega}{\mathord}{greek}{"21}                                               
\DeclareMathSymbol{\varepsilon}{\mathord}{greek}{"22}                                          
\DeclareMathSymbol{\vartheta}{\mathord}{greek}{"23}                                            
\DeclareMathSymbol{\varpi}{\mathord}{greek}{"24}                                               
\let\varrho\rho                                                                                
\let\varsigma\sigma                                                                            
\DeclareMathSymbol{\varphi}{\mathord}{greek}{"27}

\title{Автоматизация метода деконструкции Жака~Деррида на~основе техники подсказок}
\author{Савенкова~Ю.~?.}
\date{\today}

\begin{document}
\maketitle

\begin{frame}
\frametitle{Актуальность работы}
\framesubtitle{Проблематика, объект и предмет исследования}

Предложенный в~1967~г. Жаком Деррида метод деконструкции является инструментом выявления в~тексте скрытого смысла,
который позволяет очистить текст от~навязанного или тривиального содержания.
Применение техники подсказок способно автоматизировать этот процесс для~ускорения анализа текстов.

{\bf Объект исследования:} языковые модели.

{\bf Предмет исследования:} техника подсказок для~языковых моделей.
\end{frame}

\begin{frame}
\frametitle{Цель и задачи}

{\bf Цель: } автоматизировать применение метода деконструкции Жака Деррида для~выявления скрытого содержания текста
с~сохранением семантической значимости результатов.

{\bf Задачи:}

\begin{enumerate}
\item{Проанализировать этапы метода деконструкции в~наиболее формальном его~виде.}
\item{Разработать подсказки для~автоматизации каждого этапа и проведения самопроверок.}
\item{Произвести тестирование  автоматизированной реализации с~привлечением специалистов философского факультета для~оценки качества.}
\end{enumerate}
\end{frame}

\begin{frame}
\frametitle{Prompt engineering}
\framesubtitle{Что такое техника подсказок?}

Техника подсказок (Prompt engineering)~--- это процесс создания и оптимизации текстовых запросов (prompts) для моделей машинного обучения, особенно для языковых моделей, с целью получения желаемых ответов.

Основные требования к~технике подсказок:

\begin{itemize}
\item{чёткость и конкретность запроса;}
\item{Учёт контекста и специфики задачи;}
\item{использование правильных формулировок и структуры;}
\item{тестирование и оптимизация запросов для достижения лучших результатов.}
\end{itemize}
\end{frame}

\begin{frame}
\frametitle{Деконструкция}
\framesubtitle{Что предложил Жак Деррида?}

{\small

Деконструкция — это метод анализа текстов, направленный на~выявление и переосмысление бинарных оппозиций, присутствующих в~языке и культуре, а также на~раскрытие их условности и произвольности.
Основные шаги метода:

\begin{enumerate}
\item{Выявление бинарных оппозиций: определение ключевых понятий и категорий, которые противопоставляются друг другу в тексте или культурной традиции.}
\item{Анализ привилегированных терминов: изучение того, как один термин в оппозиции доминирует над другим и какие последствия это имеет для понимания текста или явления.}
\item{Переосмысление и дестабилизация оппозиций: демонстрация условности и произвольности бинарных противопоставлений, выявление их взаимозависимости и взаимопроникновения.}
\item{Создание новых интерпретаций: формирование альтернативных пониманий текста или явления на основе деконструкции бинарных оппозиций.}
\end{enumerate}

}
\end{frame}

\begin{frame}
\frametitle{Заключение}
\framesubtitle{Описание полученных результатов}

\begin{enumerate}
\item{Подобраны подсказки для~выполнения каждого этапа метода деконструкции Жака Деррида.}
\item{Полученные подсказки реализованы в~виде библиотеки с~использованием YandexGPT.}
\item{Проведено тестирование с~оценкой полученных результатов ассистентом кафедры ... Беляниным Вадимом Сергеевичем.}
\end{enumerate}
\end{frame}

\begin{frame}
\begin{center}
{\Large Спасибо за~внимание!}
\end{center}
\end{frame}
\end{document}
